\documentclass[aspectratio=169,11pt,table]{beamer}

% ================================================================
% PAQUETES
% ================================================================
\usepackage[spanish]{babel}
\usepackage[utf8]{inputenc}
\usepackage[T1]{fontenc}
\usepackage{lmodern}
\usepackage{microtype}
\usepackage{graphicx}
\usepackage[table]{xcolor}
\usepackage{booktabs}
\usepackage{array}
\usepackage{ragged2e}
\usepackage{amsmath}
\usepackage{tikz}
\usetikzlibrary{arrows.meta,positioning,calc,babel,shapes.geometric,shadows}
\usepackage[most]{tcolorbox}
\usepackage{listings}

% ================================================================
% IDENTIDAD VISUAL CCDE (replicada de la presentación de referencia)
% ================================================================
\definecolor{ccdenavy}{RGB}{10,30,64}
\definecolor{ccdeblue}{RGB}{34,53,94}
\definecolor{ccdemid}{RGB}{43,91,155}
\definecolor{ccdeaccent}{RGB}{0,119,182}
\definecolor{ccdesky}{RGB}{61,150,210}
\definecolor{ccdelight}{RGB}{115,196,223}
\definecolor{ccdeice}{RGB}{235,247,252}
\definecolor{codebg}{RGB}{245,247,250}
\definecolor{warn}{RGB}{202,93,54}
\definecolor{okgreen}{RGB}{22,128,95}

\usetheme{default}
\usefonttheme{professionalfonts}
\setbeamertemplate{navigation symbols}{}
\setbeamertemplate{section in toc}[sections numbered]
\setbeamertemplate{itemize item}{\color{ccdeaccent}\small$\blacktriangleright$}
\setbeamertemplate{itemize subitem}{\color{ccdesky}\scriptsize$\bullet$}
\setbeamertemplate{enumerate item}{\color{ccdeaccent}\insertenumlabel.}

\setbeamercolor{normal text}{fg=ccdenavy,bg=ccdeice!18}
\setbeamercolor{frametitle}{fg=white,bg=ccdenavy}
\setbeamercolor{title}{fg=white,bg=ccdenavy}
\setbeamercolor{subtitle}{fg=ccdeice}
\setbeamercolor{structure}{fg=ccdeaccent}
\setbeamercolor{block title}{fg=white,bg=ccdeblue}
\setbeamercolor{block body}{fg=ccdenavy,bg=white}

\setbeamertemplate{frametitle}{%
  \nointerlineskip
  \begin{beamercolorbox}[wd=\paperwidth,ht=0.72cm,dp=0.24cm,leftskip=0.38cm,rightskip=0.38cm]{frametitle}
    \usebeamerfont{frametitle}\insertframetitle
  \end{beamercolorbox}
}

\newcommand{\ccdelogofile}{CCDE-logo-whatsapp-cropped.png}

\setbeamertemplate{footline}{%
  \leavevmode%
  \hbox{%
    \begin{beamercolorbox}[wd=.70\paperwidth,ht=0.38cm,dp=0.12cm,leftskip=0.35cm]{author in head/foot}%
      \scriptsize\textcolor{ccdemid}{CCDE -- Árboles de decisión y métodos de ensamble}
    \end{beamercolorbox}%
    \begin{beamercolorbox}[wd=.20\paperwidth,ht=0.38cm,dp=0.12cm,center]{date in head/foot}%
      \IfFileExists{\ccdelogofile}{\includegraphics[height=0.25cm]{\ccdelogofile}}{}
    \end{beamercolorbox}%
    \begin{beamercolorbox}[wd=.10\paperwidth,ht=0.38cm,dp=0.12cm,rightskip=0.35cm]{page number in head/foot}%
      \hfill\scriptsize\textcolor{ccdeblue}{\insertframenumber/\inserttotalframenumber}
    \end{beamercolorbox}%
  }%
}

% ================================================================
% COMPONENTES
% ================================================================
\newcommand{\source}[1]{%
  \begin{tikzpicture}[remember picture,overlay]
    \node[
      anchor=south west,
      text=ccdemid,
      font=\tiny,
      align=left,
      text width=15.2cm
    ] at ([xshift=0.35cm,yshift=0.56cm]current page.south west) {Fuente: #1};
  \end{tikzpicture}%
}

\newcommand{\sourceright}[1]{%
  \begin{tikzpicture}[remember picture,overlay]
    \node[
      anchor=south west,
      text=ccdemid,
      font=\tiny,
      align=left,
      text width=6.7cm
    ] at ([xshift=8.25cm,yshift=0.56cm]current page.south west) {Fuente: #1};
  \end{tikzpicture}%
}

\newcommand{\tagpill}[1]{%
  \tikz[baseline=(x.base)]\node[
    fill=ccdeice,draw=ccdelight,rounded corners=2pt,
    inner xsep=4pt,inner ysep=2pt
  ] (x) {\scriptsize\bfseries\textcolor{ccdeblue}{#1}};%
}

\newtcolorbox{ccdebox}[1][]{
  enhanced,
  colback=white,
  colframe=ccdeblue,
  arc=1.5mm,
  boxrule=0.7pt,
  left=2.2mm,
  right=2.2mm,
  top=1.5mm,
  bottom=1.5mm,
  #1
}

\newtcolorbox{keybox}[1][]{
  enhanced,
  colback=ccdelight!13,
  colframe=ccdeaccent,
  arc=1.5mm,
  boxrule=0.8pt,
  left=2.4mm,
  right=2.4mm,
  top=1.6mm,
  bottom=1.6mm,
  fontupper=\small,
  #1
}

\newtcolorbox{contrastbox}[2]{
  enhanced,
  colback=#1!8,
  colframe=#1!80!ccdenavy,
  title=#2,
  fonttitle=\bfseries,
  coltitle=white,
  boxed title style={colback=#1!80!ccdenavy,arc=1mm,boxrule=0pt},
  attach boxed title to top left={yshift=-2mm,xshift=2mm},
  arc=1.5mm,
  boxrule=0.7pt,
  left=2.2mm,
  right=2.2mm,
  top=2.8mm,
  bottom=1.6mm
}

\lstdefinestyle{ccdepython}{
  language=Python,
  basicstyle=\ttfamily\scriptsize,
  keywordstyle=\bfseries\color{ccdeaccent},
  commentstyle=\itshape\color{okgreen},
  stringstyle=\color{warn},
  backgroundcolor=\color{codebg},
  frame=single,
  rulecolor=\color{ccdelight},
  frameround=tttt,
  showstringspaces=false,
  breaklines=true,
  columns=fullflexible,
  keepspaces=true,
  aboveskip=0cm,
  belowskip=0cm
}

\title{Árboles de decisión y métodos de ensamble}
\subtitle{De CART al bootstrap, bagging, random forests y boosting}
\author{Círculo de Ciencia de Datos y Econometría}
\date{22 de julio de 2026}

\begin{document}

% ================================================================
\begin{frame}[plain]
\begin{tikzpicture}[remember picture,overlay]
  \fill[ccdeice!35] (current page.south west) rectangle (current page.north east);
  \fill[ccdenavy] (current page.north west) rectangle ([yshift=-1.25cm]current page.north east);
  \fill[ccdeblue] ([yshift=-1.25cm]current page.north west) rectangle ([yshift=-1.52cm]current page.north east);
  \fill[ccdelight!45,opacity=0.85] ([xshift=13.6cm,yshift=-1.05cm]current page.north west) circle (1.25cm);
  \fill[ccdeaccent!18,opacity=0.7] ([xshift=0.6cm,yshift=0.5cm]current page.south west) circle (2.4cm);
  \draw[ccdelight,line width=1.4pt] ([xshift=0.75cm,yshift=-7.05cm]current page.north west) --
    ([xshift=-0.75cm,yshift=-7.05cm]current page.north east);
\end{tikzpicture}

\vspace*{0.62cm}
\begin{center}
\begin{tcolorbox}[
  enhanced,width=0.88\textwidth,colback=white,colframe=ccdelight,
  arc=2mm,boxrule=0.9pt,left=6mm,right=6mm,top=5mm,bottom=5mm,
  drop shadow={ccdenavy!22!white},borderline west={4pt}{0pt}{ccdeaccent}
]
\begin{columns}[c,onlytextwidth]
\begin{column}{0.22\textwidth}
\centering
\IfFileExists{\ccdelogofile}{\includegraphics[width=0.88\linewidth]{\ccdelogofile}}{}
\end{column}
\begin{column}{0.74\textwidth}
{\LARGE\bfseries\textcolor{ccdenavy}{Árboles de decisión y métodos de ensamble}\par}
\vspace{0.22cm}
{\large\textcolor{ccdeblue}{De CART al bootstrap, bagging, random forests y boosting}\par}
\vspace{0.38cm}
{\normalsize\bfseries\textcolor{ccdeaccent}{Círculo de Ciencia de Datos y Econometría}}\\[0.25em]
{\small 22 de julio de 2026}
\end{column}
\end{columns}
\end{tcolorbox}
\end{center}
\end{frame}

% ================================================================
\begin{frame}{Mapa de la presentación}
\begin{columns}[T,onlytextwidth]
\begin{column}{0.58\textwidth}
\tableofcontents
\end{column}
\begin{column}{0.37\textwidth}
\begin{keybox}
\textbf{Pregunta guía}\\[0.15cm]
¿Cómo pasamos de un árbol interpretable pero inestable a modelos colectivos con mejor capacidad predictiva?
\end{keybox}
\vspace{0.18cm}
\begin{ccdebox}
\small
\textbf{Fuentes troncales}\\
Géron, cap. 6; James et al., cap. 8.\\[0.12cm]
\textbf{Aplicación}\\
Código reproducible en Python.
\end{ccdebox}
\end{column}
\end{columns}
\end{frame}

% ================================================================
\section{Fundamentos de árboles}

\begin{frame}{¿Qué aprende un árbol de decisión?}
\small
\begin{columns}[c,onlytextwidth]
\begin{column}{0.45\textwidth}
\begin{itemize}
  \item Divide el espacio de predictores en \textbf{regiones rectangulares}.
  \item En cada nodo formula una regla del tipo \(x_j \leq t\).
  \item En una hoja devuelve una clase, probabilidad o promedio.
  \item Las reglas forman una secuencia legible: \textbf{si--entonces}.
\end{itemize}
\vspace{0.15cm}
\begin{keybox}
Un árbol es un modelo no paramétrico: no presupone linealidad ni una forma funcional global.
\end{keybox}
\end{column}
\begin{column}{0.52\textwidth}
\centering
\begin{tikzpicture}[
  level distance=1.28cm,
  level 1/.style={sibling distance=3.65cm},
  level 2/.style={sibling distance=1.72cm},
  every node/.style={font=\scriptsize,align=center},
  split/.style={draw,rounded corners,fill=white,text width=1.92cm,minimum height=0.58cm,drop shadow={opacity=0.10}},
  leaf/.style={draw,rounded corners,fill=ccdeice,text width=1.30cm,minimum height=0.52cm},
  edge from parent/.style={draw,-{Latex},thick,ccdemid}
]
\node[split]{ingreso \(\leq\) 2500}
  child {node[split]{deuda \(\leq\) 35\%}
    child {node[leaf,fill=ccdelight!24]{riesgo\\bajo}}
    child {node[leaf,fill=warn!12]{riesgo\\medio}}
  }
  child {node[split]{mora previa = sí}
    child {node[leaf,fill=warn!18]{riesgo\\alto}}
    child {node[leaf,fill=ccdelight!24]{riesgo\\bajo}}
  };
\node[text=ccdeaccent,font=\scriptsize\bfseries] at (-1.80,1.05) {sí};
\node[text=ccdeaccent,font=\scriptsize\bfseries] at (1.80,1.05) {no};
\end{tikzpicture}
\end{column}
\end{columns}
\source{Géron (2022), cap. 6; James et al. (2013), sec. 8.1.}
\end{frame}

% ================================================================
\begin{frame}{CART: crecer el árbol mediante divisiones binarias}
\footnotesize
\begin{columns}[T,onlytextwidth]
\begin{column}{0.52\textwidth}
\begin{contrastbox}{ccdeblue}{Búsqueda voraz}
\footnotesize
\begin{enumerate}
  \item Evaluar pares: característica \(k\) y umbral \(t_k\).
  \item Elegir la división que más reduce la impureza.
  \item Repetir el proceso en cada subconjunto.
  \item Detenerse por pureza o por restricciones.
\end{enumerate}
\end{contrastbox}
\vspace{0.22cm}
\begin{keybox}
CART no busca el árbol globalmente óptimo: toma la mejor decisión disponible en cada nodo.
\end{keybox}
\end{column}
\begin{column}{0.45\textwidth}
\begin{ccdebox}
\centering
\textbf{Función de coste de una división}
\[
J(k,t_k)=
\frac{m_{\mathrm{izq}}}{m}G_{\mathrm{izq}}
+\frac{m_{\mathrm{der}}}{m}G_{\mathrm{der}}
\]
\footnotesize
\begin{itemize}
  \item \(G\): impureza o error del nodo.
  \item \(m\): número de observaciones.
  \item Se elige \((k,t_k)\) que minimiza \(J\).
\end{itemize}
\end{ccdebox}
\vspace{0.18cm}
\centering
\tagpill{Entrenamiento: \(O(n\,m\log m)\)}\\[0.14cm]
\tagpill{Predicción: \(O(\log m)\)}
\end{column}
\end{columns}
\source{Géron (2022), ``The CART Training Algorithm'' y ``Computational Complexity''.}
\end{frame}

% ================================================================
\begin{frame}{De la raíz a la predicción}
\footnotesize
\begin{columns}[T,onlytextwidth]
\begin{column}{0.47\textwidth}
\begin{contrastbox}{ccdeblue}{Clasificación}
\footnotesize
\begin{itemize}
  \item La instancia recorre una única ruta.
  \item La clase predicha es la mayoritaria en la hoja.
  \item La probabilidad estimada es la proporción de cada clase dentro de esa hoja.
\end{itemize}
\[
\widehat{P}(Y=k\mid x)=\frac{m_{k,\mathrm{hoja}}}{m_{\mathrm{hoja}}}
\]
\end{contrastbox}
\end{column}
\begin{column}{0.47\textwidth}
\begin{contrastbox}{ccdeaccent}{Regresión}
\footnotesize
\begin{itemize}
  \item La partición del espacio es análoga.
  \item Cada hoja predice el promedio de \(y\).
  \item CART elige divisiones que reducen el error cuadrático.
\end{itemize}
\[
\widehat{y}(x)=\frac{1}{m_{\mathrm{hoja}}}
\sum_{i\in\mathrm{hoja}}y_i
\]
\end{contrastbox}
\end{column}
\end{columns}
\vspace{0.18cm}
\begin{keybox}
\centering
La frontera de decisión es escalonada: flexible ante no linealidades, pero sensible a pequeñas variaciones de los datos.
\end{keybox}
\source{Géron (2022), ``Making Predictions'', ``Estimating Class Probabilities'' y ``Regression''.}
\end{frame}

% ================================================================
\begin{frame}{¿Cómo medir la calidad de una división?}
\footnotesize
\begin{columns}[T,onlytextwidth]
\begin{column}{0.31\textwidth}
\begin{ccdebox}
\centering
\textbf{Gini}\\[-0.05cm]
\[
G=1-\sum_{k=1}^{K}p_k^2
\]
\footnotesize
Rápido, criterio predeterminado de scikit-learn. Favorece nodos puros.
\end{ccdebox}
\end{column}
\begin{column}{0.31\textwidth}
\begin{ccdebox}
\centering
\textbf{Entropía}\\[-0.05cm]
\[
H=-\sum_{k=1}^{K}p_k\log_2p_k
\]
\footnotesize
Mide desorden informacional. Suele producir árboles muy similares a Gini.
\end{ccdebox}
\end{column}
\begin{column}{0.31\textwidth}
\begin{ccdebox}
\centering
\textbf{Error de clasificación}\\[-0.05cm]
\[
E=1-\max_k p_k
\]
\footnotesize
Útil para evaluar una hoja; Gini y entropía son más sensibles al crecer el árbol.
\end{ccdebox}
\end{column}
\end{columns}
\vspace{0.28cm}
\begin{keybox}
\textbf{Lectura práctica:} la elección Gini/entropía suele importar menos que controlar profundidad, tamaño de hojas y validación fuera de muestra.
\end{keybox}
\source{Géron (2022), ``Gini Impurity or Entropy?''; James et al. (2013), sec. 8.1.2.}
\end{frame}

% ================================================================
\begin{frame}{Regularizar y podar: evitar que el árbol memorice}
\footnotesize
\begin{columns}[T,onlytextwidth]
\begin{column}{0.48\textwidth}
\begin{contrastbox}{ccdeblue}{Pre-poda}
\footnotesize
Detener el crecimiento antes de que el árbol sea demasiado complejo:
\begin{itemize}
  \item \texttt{max\_depth}
  \item \texttt{min\_samples\_split}
  \item \texttt{min\_samples\_leaf}
  \item \texttt{max\_leaf\_nodes}
\end{itemize}
\end{contrastbox}
\end{column}
\begin{column}{0.48\textwidth}
\begin{contrastbox}{ccdeaccent}{Poda por complejidad}
\footnotesize
Crecer un árbol grande y seleccionar un subárbol \(T\):
\[
\sum_{m=1}^{|T|}\sum_{i:x_i\in R_m}(y_i-\hat y_{R_m})^2+\alpha|T|
\]
\begin{itemize}
  \item \(\alpha\) penaliza hojas adicionales.
  \item Se elige con validación cruzada.
\end{itemize}
\end{contrastbox}
\end{column}
\end{columns}
\vspace{0cm}
\begin{keybox}
\centering
Más profundidad reduce sesgo, pero aumenta varianza. La complejidad debe elegirse con datos no usados para ajustar.
\end{keybox}
\source{Géron (2022), ``Regularization Hyperparameters''; James et al. (2013), ``Tree Pruning''.}
\end{frame}

% ================================================================
\begin{frame}{Ventajas y límites de un árbol individual}
\footnotesize
\begin{columns}[T,onlytextwidth]
\begin{column}{0.47\textwidth}
\begin{contrastbox}{okgreen}{Lo que hace bien}
\footnotesize
\begin{itemize}
  \item Reglas fáciles de explicar y visualizar.
  \item Captura interacciones y no linealidades con poca preparación.
  \item Sirve para clasificación y regresión.
  \item Predicción rápida.
\end{itemize}
\end{contrastbox}
\end{column}
\begin{column}{0.47\textwidth}
\begin{contrastbox}{warn}{Lo que exige cuidado}
\footnotesize
\begin{itemize}
  \item Alta varianza: sensible a cambios y a la orientación de los ejes.
  \item Un árbol profundo sobreajusta.
  \item Menor precisión que buenos ensambles.
  \item Importancia por impureza puede sesgarse.
\end{itemize}
\end{contrastbox}
\end{column}
\end{columns}
\vspace{0.08cm}
\begin{keybox}
\centering
La inestabilidad del árbol no es solo un problema: también es la oportunidad que explota \textbf{bagging}.
\end{keybox}
\source{Géron (2022), ``Sensitivity to Axis Orientation'' y ``Decision Trees Have a High Variance''; James et al. (2013), sec. 8.1.4.}
\end{frame}

% ================================================================
\section{Bootstrap y bagging}

\begin{frame}{Bootstrap: muchas muestras a partir de una sola}
\small
\begin{columns}[c,onlytextwidth]
\begin{column}{0.45\textwidth}
\begin{itemize}
  \item Partimos de una muestra \(D\) con \(n\) observaciones.
  \item Construimos \(B\) muestras \(D^{*(b)}\) de tamaño \(n\).
  \item Cada muestra se obtiene \textbf{con reemplazo}.
  \item Una observación puede repetirse o quedar fuera.
\end{itemize}
\vspace{0.16cm}
\begin{keybox}[fontupper=\footnotesize]
En promedio, cada muestra bootstrap contiene cerca del \(63.2\%\) de observaciones únicas; el resto queda \emph{out-of-bag}.
\end{keybox}
\end{column}
\begin{column}{0.51\textwidth}
\centering
\begin{tikzpicture}[
  obs/.style={circle,draw=ccdemid,fill=white,minimum size=0.52cm,font=\scriptsize},
  sample/.style={rounded corners,draw=ccdeaccent,fill=ccdeice,minimum width=4.65cm,minimum height=0.7cm,font=\scriptsize},
  >=Latex
]
\foreach \x/\v in {0/1,0.8/2,1.6/3,2.4/4,3.2/5,4/6}{
  \node[obs] (o\v) at (\x,2.25) {\v};
}
\node[font=\small\bfseries,text=ccdeblue] at (2,3.0) {Muestra original};
\node[sample] (s1) at (2,0.9) {\(\{2,5,2,1,6,3\}\)};
\node[sample] (s2) at (2,-0.15) {\(\{4,4,1,6,2,4\}\)};
\node[sample] (s3) at (2,-1.2) {\(\{3,5,6,5,1,2\}\)};
\coordinate (branch) at (-0.72,1.58);
\draw[->,thick,ccdeaccent] (2,1.98) -- (2,1.58);
\draw[thick,ccdeaccent] (2,1.58) -- (branch) -- ++(0,-2.78);
\fill[ccdeaccent] (branch) circle (1.2pt);
\draw[->,thick,ccdeaccent] (-0.72,0.9) -- (s1.west);
\draw[->,thick,ccdeaccent] (-0.72,-0.15) -- (s2.west);
\draw[->,thick,ccdeaccent] (-0.72,-1.2) -- (s3.west);
\node[anchor=west,font=\scriptsize,text=ccdemid] at (4.55,0.9) {bootstrap 1};
\node[anchor=west,font=\scriptsize,text=ccdemid] at (4.55,-0.15) {bootstrap 2};
\node[anchor=west,font=\scriptsize,text=ccdemid] at (4.55,-1.2) {bootstrap \(B\)};
\end{tikzpicture}
\end{column}
\end{columns}
\sourceright{James et al. (2013), sec. 8.2.1.}
\end{frame}

% ================================================================
\begin{frame}{Bagging: entrenar en paralelo y agregar}
\footnotesize
\begin{columns}[T,onlytextwidth]
\begin{column}{0.54\textwidth}
\centering
\begin{tikzpicture}[
  box/.style={draw,rounded corners,fill=white,text width=1.65cm,minimum height=0.58cm,align=center,font=\scriptsize,drop shadow={opacity=0.09}},
  >=Latex
]
\node[box,fill=ccdeice] (d) at (0,2.55) {datos \(D\)};
\node[box] (b1) at (-2.65,1.3) {bootstrap 1};
\node[box] (b2) at (0,1.3) {bootstrap 2};
\node[box] (b3) at (2.65,1.3) {bootstrap \(B\)};
\node[box,fill=ccdeblue!8] (t1) at (-2.65,0.0) {árbol 1};
\node[box,fill=ccdeblue!8] (t2) at (0,0.0) {árbol 2};
\node[box,fill=ccdeblue!8] (t3) at (2.65,0.0) {árbol \(B\)};
\node[box,fill=ccdelight!25,text width=2.75cm] (a) at (0,-1.40) {promedio o voto};
\foreach \x in {b1,b2,b3}{\draw[->,thick,ccdeaccent] (d) -- (\x);}
\draw[->,thick,ccdeblue] (b1) -- (t1);
\draw[->,thick,ccdeblue] (b2) -- (t2);
\draw[->,thick,ccdeblue] (b3) -- (t3);
\foreach \x in {t1,t2,t3}{\draw[->,thick,ccdeaccent] (\x) -- (a);}
\end{tikzpicture}
\end{column}
\begin{column}{0.42\textwidth}
\begin{contrastbox}{ccdeaccent}{Algoritmo de bagging}
\footnotesize
\begin{enumerate}
  \item Generar \(B\) muestras bootstrap.
  \item Ajustar un árbol profundo en cada muestra.
  \item Regresión: promediar predicciones.
  \item Clasificación: votar o promediar probabilidades.
\end{enumerate}
\end{contrastbox}
\vspace{0.16cm}
\begin{keybox}
\textbf{Bagging} = \emph{bootstrap aggregating}. Los modelos se entrenan de manera independiente y paralelizable.
\end{keybox}
\end{column}
\end{columns}
\source{Breiman (1996); James et al. (2013), sec. 8.2.1.}
\end{frame}

% ================================================================
\begin{frame}{¿Por qué funciona bagging? Reducción de varianza}
\footnotesize
\begin{columns}[T,onlytextwidth]
\begin{column}{0.51\textwidth}
\begin{ccdebox}
\footnotesize
Si \(B\) estimadores tienen varianza \(\sigma^2\) y correlación media \(\rho\), la varianza de su promedio es aproximadamente:
\[
\mathrm{Var}(\bar f)\approx
\rho\sigma^2+\frac{1-\rho}{B}\sigma^2
\]
\begin{itemize}
  \item Al crecer \(B\), desaparece la parte no correlacionada.
  \item La correlación \(\rho\) impone un piso.
  \item Árboles profundos: bajo sesgo, alta varianza.
\end{itemize}
\end{ccdebox}
\end{column}
\begin{column}{0.45\textwidth}
\begin{keybox}[fontupper=\footnotesize]
\textbf{Intuición}\\[0.12cm]
Los errores particulares de cada árbol se compensan cuando promediamos.
\end{keybox}
\vspace{0.10cm}
\begin{contrastbox}{warn}{Límite}
\footnotesize
Si todos los árboles se parecen demasiado, sus errores se mueven juntos y el promedio mejora menos.
\end{contrastbox}
\vspace{0.10cm}
\begin{contrastbox}{ccdeaccent}{Siguiente paso}
\footnotesize
Random forests sortea variables para \textbf{decorrelacionar} los árboles.
\end{contrastbox}
\end{column}
\end{columns}
\source{James et al. (2013), sec. 8.2.1.}
\end{frame}

% ================================================================
\begin{frame}{Error out-of-bag: validación ``gratis'' dentro de bagging}
\small
\begin{columns}[c,onlytextwidth]
\begin{column}{0.49\textwidth}
\begin{itemize}
  \item Para cada observación \(i\), identificar árboles que no la usaron.
  \item Predecir \(i\) solo con esos árboles OOB.
  \item Comparar la predicción agregada con \(y_i\).
  \item Promediar el error sobre todas las observaciones.
\end{itemize}
\vspace{0.18cm}
\begin{keybox}
El error OOB aproxima el error de prueba sin reservar otro conjunto ni ejecutar validación cruzada completa.
\end{keybox}
\end{column}
\begin{column}{0.47\textwidth}
\centering
\begin{tikzpicture}[
  card/.style={draw,rounded corners,minimum width=2.6cm,minimum height=0.75cm,align=center,font=\scriptsize},
  >=Latex
]
\node[card,fill=ccdeice] (x) at (0,2.05) {observación \(i\)};
\node[card,fill=warn!12] (oob) at (0,0.72) {árboles donde \(i\)\\quedó fuera};
\node[card,fill=ccdelight!24] (p) at (0,-0.72) {predicción OOB\\\(\hat y_i^{\mathrm{OOB}}\)};
\node[card,fill=white] (e) at (0,-2.05) {pérdida\\\(L(y_i,\hat y_i^{\mathrm{OOB}})\)};
\draw[->,thick,ccdeaccent] (x)--(oob);
\draw[->,thick,ccdeaccent] (oob)--(p);
\draw[->,thick,ccdeaccent] (p)--(e);
\end{tikzpicture}
\end{column}
\end{columns}
\source{James et al. (2013), ``Out-of-Bag Error Estimation''.}
\end{frame}

% ================================================================
\section{Bosques y boosting}

\begin{frame}{Random forests: bagging con árboles menos correlacionados}
\small
\begin{columns}[T,onlytextwidth]
\begin{column}{0.48\textwidth}
\begin{contrastbox}{ccdeblue}{Bagging de árboles}
\small
En cada división se consideran las \(p\) variables disponibles. Un predictor muy fuerte puede dominar muchos árboles.
\end{contrastbox}
\vspace{0.28cm}
\begin{contrastbox}{ccdeaccent}{Random forest}
\small
En cada nodo se sortea un subconjunto de \(m\) variables y la mejor división se busca solo allí.
\[
m\approx\sqrt{p}\quad\text{en clasificación}
\]
\end{contrastbox}
\end{column}
\begin{column}{0.48\textwidth}
\begin{ccdebox}
\small
\textbf{Efecto de la aleatoriedad}
\begin{itemize}
  \item Da oportunidad a predictores distintos.
  \item Reduce correlación entre árboles.
  \item Mantiene la reducción de varianza por promedio.
  \item Suele mejorar la precisión respecto a bagging puro.
\end{itemize}
\end{ccdebox}
\vspace{0.18cm}
\begin{keybox}
Bagging es el caso particular de random forest con \(m=p\).
\end{keybox}
\end{column}
\end{columns}
\source{James et al. (2013), sec. 8.2.2.}
\end{frame}

% ================================================================
\begin{frame}{Boosting: aprender de los errores anteriores}
\small
\begin{columns}[c,onlytextwidth]
\begin{column}{0.53\textwidth}
\centering
\begin{tikzpicture}[
  box/.style={draw,rounded corners,fill=white,text width=1.25cm,minimum height=0.64cm,align=center,font=\tiny,drop shadow={opacity=0.08}},
  >=Latex
]
\node[box,fill=ccdeice] (r0) at (-2.85,0.65) {datos\\\(r^{(0)}=y\)};
\node[box,fill=ccdeblue!8] (t1) at (-0.95,0.65) {árbol\\\(\hat f_1\)};
\node[box,fill=ccdeblue!8] (t2) at (0.95,0.65) {árbol sobre\\residuos};
\node[box,fill=ccdelight!25] (sum) at (2.85,0.65) {suma\\\(\sum_b\lambda\hat f_b\)};
\draw[->,thick,ccdeaccent] (r0)--(t1);
\draw[->,thick,ccdeaccent] (t1)--(t2);
\draw[->,thick,ccdeaccent] (t2)--(sum);
\node[font=\scriptsize,text=warn,align=center] at (0,-1.05) {cada árbol depende\\de los anteriores};
\draw[->,warn] (0,-0.65) -- (0,0.15);
\end{tikzpicture}
\end{column}
\begin{column}{0.43\textwidth}
\begin{contrastbox}{ccdeaccent}{Tres hiperparámetros}
\small
\begin{itemize}
  \item Número de árboles \(B\).
  \item Tasa de aprendizaje \(\lambda\).
  \item Profundidad \(d\) de cada árbol.
\end{itemize}
\end{contrastbox}
\vspace{0.18cm}
\begin{keybox}
Boosting es secuencial y no necesita muestras bootstrap: cada nuevo árbol corrige parte del error acumulado.
\end{keybox}
\end{column}
\end{columns}
\source{James et al. (2013), sec. 8.2.3, algoritmo de boosting para regresión.}
\end{frame}

% ================================================================
\begin{frame}{Árbol, bagging, random forest y boosting}
\vfill
\renewcommand{\arraystretch}{1.34}
\setlength{\tabcolsep}{5.8pt}
\scriptsize
\begin{center}
\begin{tabular}{lllll}
\toprule
\rowcolor{ccdenavy}
\parbox{1.75cm}{\textcolor{white}{\textbf{Método}}} &
\parbox{2.55cm}{\textcolor{white}{\textbf{Entrenamiento}}} &
\parbox{2.55cm}{\textcolor{white}{\textbf{Agregación}}} &
\parbox{2.60cm}{\textcolor{white}{\textbf{Fortaleza}}} &
\parbox{2.45cm}{\textcolor{white}{\textbf{Riesgo}}}\\
\midrule
\parbox{1.75cm}{\textbf{Árbol}} &
\parbox{2.55cm}{un modelo} &
\parbox{2.55cm}{una ruta} &
\parbox{2.60cm}{interpretación} &
\parbox{2.45cm}{alta varianza}\\
\rowcolor{ccdeice!50}
\parbox{1.75cm}{\textbf{Bagging}} &
\parbox{2.55cm}{bootstrap, paralelo} &
\parbox{2.55cm}{promedio o voto} &
\parbox{2.60cm}{reduce varianza} &
\parbox{2.45cm}{árboles correlacionados}\\
\parbox{1.75cm}{\textbf{Random forest}} &
\parbox{2.55cm}{bootstrap + variables aleatorias} &
\parbox{2.55cm}{promedio o voto} &
\parbox{2.60cm}{precisión robusta} &
\parbox{2.45cm}{menor interpretabilidad}\\
\rowcolor{ccdeice!50}
\parbox{1.75cm}{\textbf{Boosting}} &
\parbox{2.55cm}{secuencial, corrige errores} &
\parbox{2.55cm}{suma ponderada} &
\parbox{2.60cm}{bajo sesgo y alta precisión} &
\parbox{2.45cm}{sensible al ajuste}\\
\bottomrule
\end{tabular}
\end{center}
\vspace{0.30cm}
\begin{keybox}
\centering
No hay ganador universal: la elección depende del objetivo, costo de errores, volumen de datos y necesidad de explicación.
\end{keybox}
\vfill
\end{frame}

% ================================================================
\section{Aplicaciones reales}

\begin{frame}{¿Dónde aparecen estos modelos en la vida real?}
\footnotesize
\begin{columns}[T,onlytextwidth]
\begin{column}{0.32\textwidth}
\begin{ccdebox}
\centering
\textbf{Finanzas}\\[0.18cm]
\footnotesize
Riesgo de crédito, fraude, mora y priorización de cobranzas.
\end{ccdebox}
\vspace{0.10cm}
\begin{ccdebox}
\centering
\textbf{Salud}\\[0.18cm]
\footnotesize
Apoyo diagnóstico, riesgo de readmisión y triaje clínico.
\end{ccdebox}
\end{column}
\begin{column}{0.32\textwidth}
\begin{ccdebox}
\centering
\textbf{Marketing}\\[0.18cm]
\footnotesize
Churn, respuesta a campañas, propensión de compra y segmentación.
\end{ccdebox}
\vspace{0.10cm}
\begin{ccdebox}
\centering
\textbf{Operaciones}\\[0.18cm]
\footnotesize
Mantenimiento predictivo, fallas, demanda e inventarios.
\end{ccdebox}
\end{column}
\begin{column}{0.32\textwidth}
\begin{ccdebox}
\centering
\textbf{Sector público}\\[0.18cm]
\footnotesize
Focalización, inspección, alertas tempranas y asignación de recursos.
\end{ccdebox}
\vspace{0.10cm}
\begin{ccdebox}
\centering
\textbf{Ambiente}\\[0.18cm]
\footnotesize
Riesgo de incendios, uso de suelo, calidad del aire y agricultura.
\end{ccdebox}
\end{column}
\end{columns}
\vspace{0.14cm}
\begin{keybox}
\textbf{Principio de uso responsable:} una buena métrica predictiva no sustituye validación temporal, análisis de sesgos ni revisión humana cuando la decisión afecta derechos.
\end{keybox}
\end{frame}

% ================================================================
\begin{frame}{Mini-caso: sistema de apoyo para riesgo crediticio}
\footnotesize
\begin{columns}[T,onlytextwidth]
\begin{column}{0.45\textwidth}
\centering
\vspace{1.10cm}
\begin{tikzpicture}[
  node distance=0.46cm,
  box/.style={draw,rounded corners,fill=white,text width=4.25cm,minimum height=0.55cm,align=center,font=\small,drop shadow={opacity=0.08}},
  >=Latex
]
\node[box] (x) {ingresos, deuda, historial, estabilidad};
\node[box,fill=ccdeice,below=of x] (m) {árbol o ensamble};
\node[box,fill=ccdelight!22,below=of m] (p) {probabilidad de mora};
\node[box,below=of p] (d) {revisión + decisión + monitoreo};
\foreach \a/\b in {x/m,m/p,p/d}{\draw[->,thick,ccdeaccent] (\a)--(\b);}
\end{tikzpicture}
\end{column}
\begin{column}{0.51\textwidth}
\begin{keybox}
\textbf{Pregunta:} ¿quién requiere una evaluación crediticia más cuidadosa?
\end{keybox}
\vspace{0.15cm}
\footnotesize
\begin{itemize}
  \item \textbf{Métrica:} ROC-AUC, recall de mora y costo esperado.
  \item \textbf{Validación:} corte temporal, no solo partición aleatoria.
  \item \textbf{Explicación:} reglas, importancia por permutación o SHAP.
  \item \textbf{Riesgo:} proxies de atributos protegidos y deriva económica.
  \item \textbf{Uso:} apoyo a la decisión, con mecanismos de revisión.
\end{itemize}
\end{column}
\end{columns}
\end{frame}

% ================================================================
\section{Código Python}

\begin{frame}[fragile]{Python 1/3: datos, partición y árbol regularizado}
\scriptsize
\begin{columns}[T,onlytextwidth]
\begin{column}{0.62\textwidth}
\begin{lstlisting}[style=ccdepython]
from sklearn.datasets import load_breast_cancer
from sklearn.model_selection import train_test_split
from sklearn.tree import DecisionTreeClassifier
data = load_breast_cancer(as_frame=True)
X_train, X_test, y_train, y_test = train_test_split(
    data.data, data.target, test_size=0.25,
    stratify=data.target, random_state=42
)
tree = DecisionTreeClassifier(
    max_depth=3, min_samples_leaf=8,
    random_state=42
)
tree.fit(X_train, y_train)
\end{lstlisting}
\end{column}
\begin{column}{0.34\textwidth}
\begin{keybox}
\textbf{Decisiones metodológicas}
\begin{itemize}
  \item \texttt{stratify}: conserva proporciones.
  \item Profundidad y hoja mínima: regularizan.
  \item \texttt{random\_state}: reproducibilidad.
\end{itemize}
\end{keybox}
\vspace{0cm}
\begin{ccdebox}
\scriptsize
Datos: Breast Cancer Wisconsin, incluidos en scikit-learn. El objetivo es clasificación binaria.
\end{ccdebox}
\end{column}
\end{columns}
\end{frame}

% ================================================================
\begin{frame}[fragile]{Python 2/3: bagging y random forest}
\begin{columns}[T,onlytextwidth]
\begin{column}{0.62\textwidth}
\begin{lstlisting}[style=ccdepython]
from sklearn.ensemble import (
    BaggingClassifier, RandomForestClassifier
)

bag = BaggingClassifier(
    estimator=DecisionTreeClassifier(random_state=42),
    n_estimators=300, max_samples=0.8,
    bootstrap=True, oob_score=True,
    n_jobs=-1, random_state=42
)

rf = RandomForestClassifier(
    n_estimators=500, max_features="sqrt",
    min_samples_leaf=2, oob_score=True,
    n_jobs=-1, random_state=42
)

bag.fit(X_train, y_train)
rf.fit(X_train, y_train)
\end{lstlisting}
\end{column}
\begin{column}{0.34\textwidth}
\begin{keybox}
\textbf{Correspondencia teórica}
\begin{itemize}
  \item \texttt{bootstrap=True}: remuestreo.
  \item \texttt{oob\_score=True}: error OOB.
  \item \texttt{max\_features}: decorrelación.
  \item \texttt{n\_jobs=-1}: paralelismo.
\end{itemize}
\end{keybox}
\end{column}
\end{columns}
\end{frame}

% ================================================================
\begin{frame}[fragile]{Python 3/3: boosting y evaluación}
\begin{columns}[T,onlytextwidth]
\begin{column}{0.61\textwidth}
\begin{lstlisting}[style=ccdepython]
from sklearn.ensemble import GradientBoostingClassifier
from sklearn.metrics import accuracy_score, roc_auc_score

boost = GradientBoostingClassifier(
    n_estimators=150, learning_rate=0.05,
    max_depth=2, random_state=42
)
boost.fit(X_train, y_train)

for name, model in models.items():
    pred = model.predict(X_test)
    prob = model.predict_proba(X_test)[:, 1]
    print(name,
          accuracy_score(y_test, pred),
          roc_auc_score(y_test, prob))
\end{lstlisting}
\end{column}
\begin{column}{0.35\textwidth}
\begin{keybox}
\textbf{No evaluar solo accuracy}
\begin{itemize}
  \item ROC-AUC: ordenamiento probabilístico.
  \item Matriz de confusión: tipos de error.
  \item Calibración: calidad de probabilidades.
  \item Métrica de negocio: costo real.
\end{itemize}
\end{keybox}
\end{column}
\end{columns}
\end{frame}

% ================================================================
\begin{frame}{Resultado reproducible del ejemplo}
\begin{columns}[T,onlytextwidth]
\begin{column}{0.58\textwidth}
\vspace{2.20cm}
\renewcommand{\arraystretch}{1.32}
\footnotesize
\begin{tabular}{lccc}
\toprule
\rowcolor{ccdenavy}
\textcolor{white}{\textbf{Modelo}} &
\textcolor{white}{\textbf{Accuracy}} &
\textcolor{white}{\textbf{ROC-AUC}} &
\textcolor{white}{\textbf{OOB}}\\
\midrule
Árbol & 0.951 & 0.962 & --\\
\rowcolor{ccdeice!50}
Bagging & 0.951 & 0.993 & 0.960\\
Random forest & \textbf{0.958} & \textbf{0.994} & 0.955\\
\rowcolor{ccdeice!50}
Gradient boosting & 0.951 & 0.993 & --\\
\bottomrule
\end{tabular}
\end{column}
\begin{column}{0.38\textwidth}
\begin{keybox}
\textbf{Lectura}
\begin{itemize}
  \item Los ensambles elevan claramente ROC-AUC.
  \item OOB es cercano al desempeño de prueba.
  \item Diferencias pequeñas requieren repetición o validación cruzada.
\end{itemize}
\end{keybox}
\vspace{0cm}
\begin{ccdebox}
\footnotesize
Una sola partición ilustra el flujo, pero no basta para afirmar superioridad general.
\end{ccdebox}
\end{column}
\end{columns}
\vspace{0cm}
\source{Ejecución local de \texttt{codigo\_demo\_arboles\_ensambles.py}, scikit-learn 1.8.0, semilla 42.}
\end{frame}

% ================================================================
\begin{frame}{Interpretar y validar antes de desplegar}
\begin{columns}[T,onlytextwidth]
\begin{column}{0.32\textwidth}
\begin{ccdebox}
\centering
\textbf{Desempeño}\\[0.16cm]
\small
Validación cruzada o temporal, métricas por subgrupo, calibración y curvas de aprendizaje.
\end{ccdebox}
\end{column}
\begin{column}{0.32\textwidth}
\begin{ccdebox}
\centering
\textbf{Interpretación}\\[0.16cm]
\small
Reglas del árbol, importancia por permutación, dependencia parcial y explicaciones locales.
\end{ccdebox}
\end{column}
\begin{column}{0.32\textwidth}
\begin{ccdebox}
\centering
\textbf{Operación}\\[0.16cm]
\small
Latencia, monitoreo de deriva, umbrales, trazabilidad y proceso de revisión humana.
\end{ccdebox}
\end{column}
\end{columns}
\vspace{0.30cm}
\begin{keybox}
\centering
\textbf{Evitar fuga de información:} toda transformación, selección y ajuste debe aprenderse solo con los datos de entrenamiento.
\end{keybox}
\end{frame}

% ================================================================
\section{Cierre y fuentes}

\begin{frame}{Ideas que deben quedar}
\begin{columns}[T,onlytextwidth]
\begin{column}{0.58\textwidth}
\begin{enumerate}
  \item Un árbol convierte datos en reglas y regiones fáciles de explicar.
  \item CART crece de manera voraz; regularización y poda controlan el sobreajuste.
  \item El bootstrap permite simular nuevas muestras desde los datos observados.
  \item Bagging reduce varianza; random forests además decorrelaciona árboles.
  \item Boosting construye un modelo secuencial que corrige errores.
  \item La elección final depende de validación, costo de errores y contexto de uso.
\end{enumerate}
\end{column}
\begin{column}{0.37\textwidth}
\begin{keybox}
\centering
\large
\textbf{Un árbol explica.}\\[0.18cm]
\textbf{Un bosque estabiliza.}\\[0.18cm]
\textbf{Boosting refina.}
\end{keybox}
\vspace{0.24cm}
\centering
\tagpill{interpretabilidad}\\[0.12cm]
\tagpill{varianza}\\[0.12cm]
\tagpill{generalización}
\end{column}
\end{columns}
\end{frame}

% ================================================================
\begin{frame}{Bibliografía y alcance de capítulos}
\small
\begin{ccdebox}
\textbf{Géron, A.} (2022). \emph{Hands-On Machine Learning with Scikit-Learn, Keras, and TensorFlow}, 3.ª ed., O'Reilly.\\
\textbf{Capítulo 6:} entrenamiento y visualización; predicción; probabilidades; CART; complejidad; Gini/entropía; regularización; regresión; sensibilidad a ejes; alta varianza.
\end{ccdebox}
\vspace{0.12cm}
\begin{ccdebox}
\textbf{James, G.; Witten, D.; Hastie, T.; Tibshirani, R.} (2013). \emph{An Introduction to Statistical Learning with Applications in R}. Springer.\\
\textbf{Capítulo 8:} árboles de regresión y clasificación; árboles vs. modelos lineales; ventajas/desventajas; bagging; OOB; importancia; random forests; boosting; laboratorio.
\end{ccdebox}
\vspace{0.12cm}
\begin{keybox}
Los ejercicios del final de capítulo no se reproducen; el laboratorio de ISLR se traduce conceptualmente de R a Python para cumplir el objetivo práctico.
\end{keybox}
\end{frame}

% ================================================================
\begin{frame}{Referencias complementarias}
\small
\begin{columns}[T,onlytextwidth]
\begin{column}{0.48\textwidth}
\begin{ccdebox}
\textbf{Breiman, L.} (1996). ``Bagging Predictors''. \emph{Machine Learning}, 24, 123--140.
\end{ccdebox}
\vspace{0.18cm}
\begin{ccdebox}
\textbf{Breiman, L.} (2001). ``Random Forests''. \emph{Machine Learning}, 45, 5--32.
\end{ccdebox}
\end{column}
\begin{column}{0.48\textwidth}
\begin{ccdebox}
\textbf{scikit-learn developers.} Documentación de \texttt{tree} y \texttt{ensemble}: \texttt{scikit-learn.org}.
\end{ccdebox}
\vspace{0.18cm}
\begin{ccdebox}
\textbf{Código asociado:}\\
\texttt{codigo\_demo\_arboles\_ensambles.py}
\end{ccdebox}
\end{column}
\end{columns}
\vfill
\begin{keybox}
\centering
Presentación elaborada a partir de los dos capítulos indicados y con el sistema visual CCDE de la carpeta de referencia.
\end{keybox}
\end{frame}

\end{document}
